\chapter{Antecedentes y Alcance}

\section{Contexto del Proyecto}

El sistema eléctrico colombiano presenta una alta dependencia de la generación hidráulica, históricamente entre el 70\% y el 80\% de la energía producida. Esta concentración expone al sistema a riesgos asociados a la variabilidad climática. En ese marco, el Estado ha promovido las Fuentes No Convencionales de Energía Renovable (FNCER) mediante instrumentos como:

\begin{itemize}
    \item Ley 1715 de 2014 y sus actualizaciones
    \item Resolución CREG 030 de 2018 (conexión de AGPE/AGGE)
    \item Resolución CREG 174 de 2021 (lineamientos de autogeneración a pequeña escala)
    \item Resolución CREG 025 de 1995 -- Código de Redes (y modificaciones)
\end{itemize}

\section{Alcance del Estudio}

El supermercado MAS X MENOS sede Guarín adelanta el proyecto \textbf{SOLAR MAS X MENOS GUARÍN} (141,6~kWp). Por su potencia AC (120~kW), se clasifica como autogenerador a pequeña escala (AGPE) y requiere estudio de conexión simplificado ante ESSA.

El alcance técnico del presente estudio cubre:

\begin{enumerate}
    \item Caracterización del circuito 10-502 (SE Conucos) con base en datos oficiales del OR.
    \item Modelado del punto de conexión en el nodo 3272966 y del transformador de 150~kVA.
    \item Análisis de flujo de carga en operación normal (9:00, 12:00 y 15:00).
    \item Análisis de cortocircuito trifásico y monofásico (IEC~60909).
    \item Verificación de protecciones del circuito y del generador (Acuerdo CNO 1322).
    \item Estimación de pérdidas técnicas con y sin SSFV.
\end{enumerate}

\section{Objetivo General}

Analizar la viabilidad técnica de la conexión del proyecto fotovoltaico de 141,6~kWp en el nodo 3272966 del circuito 10-502 de ESSA (13,8~kV), determinando los impactos sobre tensiones, cargabilidades, cortocircuitos, protecciones y pérdidas.

\section{Objetivos Específicos}

\begin{enumerate}
    \item Caracterizar la red de influencia con los parámetros eléctricos de 215 tramos y 51 cargas del circuito.
    \item Evaluar tensiones y cargabilidades con/sin generación en las franjas horarias críticas.
    \item Calcular corrientes de cortocircuito trifásico y monofásico en el POC y nodos aledaños.
    \item Contrastar los ajustes del relé MICOM P142 del circuito 10-502 frente al aporte del SSFV y elaborar la coordinación de protecciones (curva TCC).
    \item Cuantificar la variación de pérdidas técnicas atribuible a la generación distribuida.
    \item Emitir concepto técnico de conexión y recomendaciones de seguimiento.
\end{enumerate}
